% Generated by Listas/make_statements.py from pratica-lecture-02.md, item 12. Do not edit.
\textbf{Source.} Tadelis (2013), \emph{Game Theory: An Introduction}, Exercise 5.8 ``The \(n\)-Firm Cournot Model'', pp.~96--97 {[}PDF p.~115--116{]}. Answer in the source: none (the preface, p.~xiii, says about half of the exercises are solved on Princeton UP's website). Reproduced verbatim (the printed ``\(Q = q_i + \cdots + q_n\)'' is kept as printed).



\begin{quote}

\textbf{The \(n\)-Firm Cournot Model:} Suppose there are \(n\) firms in the Cournot oligopoly model. Let \(q_i\) denote the quantity produced by firm \(i\), and let \(Q = q_i + \cdots + q_n\) denote the aggregate production. Let \(P(Q)\) denote the market clearing price (when demand equals \(Q\)) and assume that the inverse demand function is given by \(P(Q) = a - Q\), where \(Q \le a\). Assume that firms have no fixed cost and that the cost of producing quantity \(q_i\) is \(cq_i\) (all firms have the same marginal cost, and assume that \(c < a\)).



\begin{enumerate}

\def\labelenumi{\alph{enumi}.}

\item

  Model this as a normal-form game.

\item

  What is the Nash (Cournot) equilibrium of the game in which firms choose their quantities simultaneously?

\item

  What happens to the equilibrium price as \(n\) approaches infinity? Is this familiar?

\end{enumerate}

\end{quote}

